% id: 95-04
% book: 베이직쎈 확률과 통계 · 05 확률변수와 확률분포
% section: 유형 056 이산확률변수의 평균, 분산, 표준편차; 확률분포가 주어진 경우
% figure: none
% answer: 
% answer_source: 
% variant_level: 0 (원본 전사)
% 이 파일은 build.mjs 가 items.json 에서 만든다. 고칠 때는 items.json 을 고친다.
\smash{\raisebox{6.5mm}{\dmkichul{베이직쎈 확률과 통계 95쪽 04번}}}
확률변수 $X$의 확률질량함수가\par\smallskip\dispeq{$\mathrm{P}(X=x)=\begin{cases} k & (x=1{,}\mkern3mu\mathopen{}3)\\[4pt] \dfrac{1}{6} & (x=2{,}\mkern3mu\mathopen{}4)\end{cases}$}\par\smallskip\noindent 일 때, $X$의 분산은? \dmcondfill{$k$는 상수이다.}{}
\dmchoicesauto{\rule[-3ex]{0pt}{8ex}}{$\dfrac{1}{3}$}{$\dfrac{5}{9}$}{$\dfrac{7}{9}$}{$1$}{$\dfrac{11}{9}$}
